\section{从小块求和到定积分}
\label{sec:02-Calculus/04-Integration/01}

你的行车记录仪每隔半分钟记一次速度。从 0 秒到 3 分钟，读数依次是（单位：米/秒）：

\[
0,\;8,\;14,\;18,\;20,\;18,\;12.
\]

这段时间里车走了多远？

如果速度恒定，答案就是小学乘法——距离等于速度乘时间。可这七个读数没一个相同：起步时是零，拉到 20 又降回 12。每个读数只在它被记录的那一瞬间''对''——它不负责描述两个读数之间的空白区。

手头的工具只有一个恒速公式和一串离散读数。两个东西不匹配。怎么办？

\subsection{没有恒速，就把每一小段当作恒速来近似}

一个直接的想法：既然每半分钟读一次数，不妨假设在这半分钟里速度大致没变——就用读数代表整段的平均速度。第一段（0 到 0.5 分钟，即 30 秒）：速度大约 0 米/秒，距离约 \(0\times30=0\) 米。第二段：8 米/秒跑 30 秒，240 米。继续下去：

\[
\begin{aligned}
\text{近似总距离} &= 0\times30 + 8\times30 + 14\times30 + 18\times30 + 20\times30 + 18\times30 + 12\times30 \\
&= (0+8+14+18+20+18+12)\times30 \\
&= 90\times30 = 2700\text{ 米}.
\end{aligned}
\]

这个数有多可靠？半分钟之内速度显然会变——踩油门时速度是连续滑上去的，不是跳变。但半分钟内的变化幅度不算大，所以 2700 应该离真相不远。如果记录间隔缩短到每秒一次，每段内的变化就更小，近似就更可靠——代价是要求和的小块数从 7 个暴涨到 180 个。计算量打了折扣换精度。

这个思路的骨架可以抽象到任意连续函数：把区间 \([a,b]\) 切成小段，在第 \(i\) 段里挑一个代表点 \(x_i^*\)，用 \(f(x_i^*)\) 当作这一整段上函数值的统一代表，乘以段宽 \(\Delta x_i\)，全部求和。切得越细，小段内的波动越小，用单点值代替整段的误差就越小。

\subsection{切分、选点、求和：Riemann 和的三个自由度}

把 \([a,b]\) 切成 \(n\) 段：

\[
a=x_0<x_1<x_2<\cdots<x_{n-1}<x_n=b,
\]

每段宽度 \(\Delta x_i=x_i-x_{i-1}\)。在第 \(i\) 段 \([x_{i-1},x_i]\) 内任意挑一个样本点 \(x_i^*\)。然后求和：

\[
\sum_{i=1}^{n} f(x_i^*)\,\Delta x_i.
\]

这就是 Riemann 和。它有三个自由选择项：怎么切（均匀还是不均匀）、切多宽、每段样本取在哪里。在极限到来之前，不同选择给出的和数值可以差不少。

用 \(f(x)=x\) 在 \([0,1]\) 上检验。等分成 \(n\) 段，\(\Delta x=1/n\)。如果取右端点，第 \(i\) 段样本在 \(i/n\)：

\[
S_n^{\text{右}} = \sum_{i=1}^{n} \frac{i}{n}\cdot\frac{1}{n}
= \frac{1}{n^2}\cdot\frac{n(n+1)}{2}
= \frac{n+1}{2n} \to \frac12.
\]

取左端点呢？第 \(i\) 段样本在 \((i-1)/n\)：

\[
S_n^{\text{左}} = \sum_{i=1}^{n} \frac{i-1}{n}\cdot\frac{1}{n}
= \frac{1}{n^2}\cdot\frac{n(n-1)}{2}
= \frac{n-1}{2n} \to \frac12.
\]

有限 \(n\) 时左右差着 \(1/n\)，但极限相同——都是 \(1/2\)。图形上这就是 \([0,1]\) 上直角三角形的面积。但请注意：这里你没用''三角形面积等于 \(\frac12\times\text{底}\times\text{高}\)''这条几何公式——你从求和极限直接推到了 \(1/2\)。这让积分为更复杂的累积量（而非只是一些简单几何形状的面积）提供了统一定义的可能性。

\subsection{极限真的与选择无关吗：Dirichlet 函数的一瓢冷水}

上面两套样本点选了不同的位置，极限都收敛到同一个数。这是巧合还是必然？如果被积函数''行为不太好''，样本点的选择会不会导致截然不同的极限？

考虑 Dirichlet 函数在 \([0,1]\) 上：

\[
f(x)=\begin{cases}
1, & x\text{ 为有理数},\\
0, & x\text{ 为无理数}.
\end{cases}
\]

任何小区间——无论多窄——里面同时住着有理数和无理数。如果你在每个子区间里全部选有理数作为样本点，Riemann 和恒为 \(1\times(1-0)=1\)。全部选无理数，和恒为 \(0\)。不管分割做得多细，这两个数之间的鸿沟纹丝不动——极限不唯一，函数不可积。

这个反例把''图像看上去围成一块面积''的直觉连根拔起。积分的存在不是一个视觉事实，而是一个极限的唯一性事实。Riemann 可积的准确定义是：当最大子区间宽度 \(\max\Delta x_i\to 0\) 时，无论你怎样选择分割和样本点，Riemann 和都收敛到同一个数——这个数就是 \(\int_a^b f(x)\,dx\)。

好消息是：闭区间上的连续函数必定 Riemann 可积。连续性保住了子区间足够细时函数值波动足够小——不同样本点之间的分歧被压到可以忽略。更一般的可积条件也允许有限个间断点，只要你切分时能把它们隔离在宽度可以任意缩小的区间里。

\(dx\) 在积分号里不仅标识积分变量。它记录了累加时每个小区间的''微小宽度''，同时承载了量纲信息：\(f\) 的单位是米/秒，\(dx\) 是秒，积出来就是米；密度千克/米乘 \(dx\)（米）得到质量千克。

\subsection{定积分是带符号的净累积}

如果 \(f(x)>0\)，小矩形贡献为正；\(f(x)<0\)，贡献为负。因此

\[
\int_a^b f(x)\,dx
\]

表示的是带符号的净累积，不是曲线与横轴之间涂了颜色的总面积。

速度函数 \(v(t)=2t-4\) 在 \([0,4]\) 上的积分：前两秒速度为负（向后运动），后两秒为正（向前）。分段算——\(\int_0^2(2t-4)\,dt\) 得 \(-4\)，\(\int_2^4(2t-4)\,dt\) 得 \(4\)，净积分为零（回到起点）。但如果是总路程——管方向，只计里程——应该积 \(|v(t)|\)，结果为 \(4+4=8\)。

这一区别在物理上以各种面目反复出现：净位移 vs 总路程，净流量 vs 总吞吐量，净电荷转移 vs 总电荷移动量。认清楚你要的是带符号量还是绝对值量，选对被积函数，比积分的计算细节更先出错。

积分还有两个基本操作属性：

\[
\int_b^a f(x)\,dx = -\int_a^b f(x)\,dx,\qquad
\int_a^c f(x)\,dx = \int_a^b f(x)\,dx + \int_b^c f(x)\,dx.
\]

反向积分的变号来自''有向宽度''的理解：\(dx\) 为负方向时每个小宽度反转，总和自然反号。区间可加性来自把所有小块重新编组——先按 \(a\to b\) 和 \(b\to c\) 两段分别累加再合并，结果与从 \(a\) 一口气加到 \(c\) 完全一致。

\subsection{从离散读数到连续极限}

行车记录仪把连续的速度曲线离散成了七个点，你用常速近似把七个点拼成了一个估计距离。Riemann 和的极限过程做的是同一件事，只是把离散化的密度推到无限——切碎、采样、求和、取极限。最后的极限值就是定积分，一个不依赖具体切割和采样方式的确定数值。

定积分的存在性和唯一性是保证，但直接走 Riemann 和极限这条路算积分，对绝大多数函数来说都苦难重重。下一节要揭开的，是整个微积分中最惊人的一个事实：求这个极限，居然可以通过''找导数的逆过程''来一笔勾销。
